\subsection{Conditional trajectories of outcome-binary mean flows}
\label{app:theory-objective-inertness}

Under Assumption~\ref{ass:categorical-mean-flow}, every outcome-binary
estimator with a positive correctness coefficient induces the same
conditional trajectory in $q$. Its advantage changes only the rate at which
that trajectory is traversed. Conditioning on correctness averages over all
correct modes, so a reward-dependent advantage cannot distinguish their
identities.

\paragraph{Proof of Theorem~\ref{thm:objective-inertness}.}

\begin{proof}
\enforcehalffinalproseline
The score identities used in Lemma~\ref{lem:grpo-mean} give
$\E[\nabla_z\log p_{Z_i}\mid R_i=1]=\nabla_z\log P$ and
$\E[\nabla_z\log p_{Z_i}\mid R_i=0]=\nabla_z\log(1-P)$.
Independence of the draws means that conditioning also on the other rows'
rewards does not change either conditional mean. Since $a(R_i,R)$ depends
only on the binary rewards, conditioning on their complete vector and
summing the $R$ successful and $G-R$ failed rows gives
\[
  \E[\widehat g_a]
  =\frac1G\Big(\E[R\,a(1,R)]\nabla_z\log P
   +\E[(G-R)\,a(0,R)]\nabla_z\log(1-P)\Big).
\]
Substituting $\nabla_z\log P=\nabla_zP/P$ and
$\nabla_z\log(1-P)=-\nabla_zP/(1-P)$ yields
Equation~\eqref{eq:objective-inertness}. Hence
$\frac{d}{dt}\log(q_c/q_d)=\dot z_c-\dot z_d=c_a(P)P(1-P)(q_c-q_d)$.
For $c_a>0$, setting $d\tau=c_a(P)P(1-P)\,dt$ gives
Equation~\eqref{eq:grpo-replicator}.
\end{proof}

For a fixed reward-dependent advantage, the binomial identities
$R\binom GR=G\binom{G-1}{R-1}$ and
$(G-R)\binom GR=G\binom{G-1}{R}$ give
\begin{equation}
 c_a(P)=\sum_{j=0}^{G-1}\binom{G-1}{j}P^j(1-P)^{G-1-j}
 \big[a(1,j+1)-a(0,j)\big],
 \label{eq:inertness-polynomial}
\end{equation}
a polynomial in $P$ of degree at most $G-1$. Finite advantage values
therefore give a smooth, bounded coefficient on $[0,1]$.

The coefficients in this Bernstein expansion are
\begin{equation}
 d_j:=a(1,j+1)-a(0,j),\qquad j=0,\ldots,G-1 .
 \label{eq:leave-one-out-gap}
\end{equation}
With the other $G-1$ rows fixed at $j$ successes, $d_j$ is the change in a
row's advantage when its reward changes from failure to success. The
Bernstein basis is nonnegative on $[0,1]$ and sums to one, so $c_a(P)$ lies
between $\min_jd_j$ and $\max_jd_j$. Each basis function is strictly positive
on $(0,1)$; hence nonnegative gaps, with at least one positive, imply
$c_a(P)>0$ throughout the interior. This condition is sufficient:
mixed-sign gaps alone do not determine the sign of $c_a$.

For Dr.GRPO, every $d_j=(G-1)/G$, giving the constant coefficient in
Lemma~\ref{lem:grpo-mean}. For the centered MaxRL estimator,
$d_0=G-1$ and $d_j=G/(j+1)$ for $1\le j\le G-1$.
In particular, $a(1,G)=0$ and $a(0,G-1)=-1$ give $d_{G-1}=1$.
Substitution into Equation~\eqref{eq:inertness-polynomial} gives
$[1-(1-P)^{G-1}]/P$, as in Lemma~\ref{lem:maxrl-mean}.

\begin{corollary}[Collapse for correctness-improving outcome-binary mean flows]
\enforcehalffinalproseline
\label{cor:class-collapse}
Under Assumption~\ref{ass:categorical-mean-flow}, let $a(R_i,R)$ be a fixed,
bounded advantage depending only on the binary rewards. Suppose the gaps
in Equation~\eqref{eq:leave-one-out-gap} satisfy $d_j\ge0$ for every $j$,
with at least one strictly positive gap. If $q(0)$ has a unique largest
coordinate $c^\star$, then
\[
 P(t)\longrightarrow1,\qquad
 q_{c^\star}(t)\longrightarrow1,\qquad
 q_c(t)\longrightarrow0\quad(c\ne c^\star).
\]
The surviving correct mode is the one with the largest initial probability;
the choice of $a$ changes the rate along the conditional trajectory.
\end{corollary}

\begin{proof}
\enforcehalffinalproseline
By Theorem~\ref{thm:objective-inertness}, the flow is
$\dot z=c_a(P)\nabla_zP$. Equation~\eqref{eq:inertness-polynomial} makes
$c_a$ smooth and bounded on $[0,1]$. If $d_{j_0}>0$, the nonnegative-gap
hypothesis gives
\[
 c_a(P)\ge\binom{G-1}{j_0}P^{j_0}(1-P)^{G-1-j_0}d_{j_0}>0
 \qquad(0<P<1).
\]
These are the coefficient properties used in the proof of
Theorem~\ref{thm:grpo-collapse}: the logit flow exists for all finite times,
$P\to1$, and the effective time $\tau\to\infty$. The conditional
replicator equation~\eqref{eq:grpo-replicator} then eliminates all coordinates
strictly below the initial maximum. It contains neither the advantage
coefficients nor the rollout group size, so the surviving mode depends only
on $q(0)$.
\end{proof}

The nonnegative-gap condition is sufficient rather than necessary. An
estimator with mixed-sign gaps also has the same limiting behavior if
$c_a(P)>0$ throughout $(0,1)$. An interior zero of $c_a$ is a stationary
correctness level for the mean flow, while a negative coefficient gives
$\dot P=c_a(P)\lVert\nabla_zP\rVert_2^2<0$ at an interior policy.
The collapse conclusion requires positivity and an initial imbalance under
Assumption~\ref{ass:categorical-mean-flow}; it does not describe shared-prompt
or sampled neural updates.

\begin{remark}[Group size in the deterministic conditional flow]
\label{rem:group-size-inert}
For each $G\ge2$, group size enters the mean flow through $c_a(P)$.
When this coefficient is positive, the time change in
Theorem~\ref{thm:objective-inertness} gives the same conditional replicator
equation. Increasing $G$ can change the speed of the mean flow but cannot
change the surviving mode under the conditions of
Corollary~\ref{cor:class-collapse}.
Group size also changes the probability of a mixed-reward group in
Lemma~\ref{lem:replay-gradient-availability}, which is small near both
$P=0$ and $P=1$, and the MaxRL coefficient in
Corollary~\ref{cor:maxrl-starvation}. These properties concern total
correctness $P$; they do not provide a probability floor for an individual
correct mode. Gradient covariance and finite-sample symmetry breaking can depend on $G$,
so sampled updates need not share this equivalence.
\end{remark}

\begin{corollary}[MaxRL amplification of the correctness gradient]
\enforcehalffinalproseline
\label{cor:maxrl-starvation}
For the coefficients of Lemmas~\ref{lem:grpo-mean} and~\ref{lem:maxrl-mean},
\[
 \frac{c_G^{\mathrm{MaxRL}}(P)}{c_G^{\mathrm{Dr.GRPO}}}
 =\frac{G\big[1-(1-P)^{G-1}\big]}{(G-1)P}.
\]
This ratio is nonincreasing in $P$, with limits $G$ as $P\to0$ and
$G/(G-1)$ as $P\to1$. It is strictly decreasing for $G>2$ and constant
at $2$ for $G=2$. At the same policy, MaxRL therefore multiplies the
Dr.GRPO correctness-gradient coefficient by a factor approaching $G$ when
correctness is small. The probability of a nonzero fresh task component
remains bounded by the mixed-group probability in
Lemma~\ref{lem:replay-gradient-availability} for both estimators.
\end{corollary}

\begin{proof}
\enforcehalffinalproseline
Divide $c_G^{\mathrm{MaxRL}}(P)=[1-(1-P)^{G-1}]/P$ by $(G-1)/G$.
The finite-series identity
$c_G^{\mathrm{MaxRL}}(P)=\sum_{j=0}^{G-2}(1-P)^j$
gives the endpoint limits and monotonicity. For $G=2$ the sum is the
constant $1$; for $G>2$ it contains the strictly decreasing term $1-P$.
\end{proof}

